En 2025, Jean a acheté un téléphone portable d'une valeur de \np{1060}~\euro.

On suppose que la valeur du téléphone diminue de 100~\euro{} chaque année.

Soit $\left(u_n\right)$ la suite modélisant la valeur du téléphone, en euros, l'année $2025+n$ pour tout entier naturel $n$.

Ainsi $u_0 = \np{1060}$.

\begin{enumerate}
	\item Calculer $u_1$ et $u_2$. Interpréter la valeur de $u_2$ dans le contexte de l'exercice.
	\item Exprimer $u_{n+1}$ en fonction de $u_n$ pour tout entier naturel $n$.
	\item En déduire la nature de la suite $\left(u_n\right)$ et préciser la raison de la suite.
	\item Exprimer $u_n$ en fonction de $n$.
	\item À partir de quelle année la valeur du téléphone aura-t-elle perdu plus de la moitié de sa valeur initiale ?
\end{enumerate}
